\documentclass[11pt]{article}
\usepackage[utf8]{inputenc}
\usepackage{geometry}
\geometry{margin=1in}
\usepackage{amsmath}
\usepackage[T1]{fontenc}
\usepackage{lmodern}

\title{A Comparative Analysis of Feature Selection Methods:\\
Information Imbalance, Differentiable Information Imbalance,\\
and Mutual Information}
\author{Salcuni Eleonora}
\date{April 2026}

\begin{document}

\maketitle

\section*{Abstract}

Feature selection is a fundamental challenge in machine learning and
statistics, aiming to identify the most informative variables in
high-dimensional datasets to build interpretable and generalisable models.
Mutual Information (MI) measures statistical dependence between each
feature and the target independently; through greedy conditional
selection (JMI) it can account for already-selected features.
Information Imbalance (II), based on neighbourhood rank structure,
evaluates features per-feature or jointly through backward elimination
over the full feature space.
Differentiable Information Imbalance (DII) extends II through
simultaneous optimisation of continuous feature weights via gradient
descent.
A controlled simulation on a synthetic dataset of 27 features
($N=2000$) with known ground truth compares methods under two
settings: single-feature and joint features comparison.
Performance is measured by Kendall's $\tau$ with the true
importance ranking (higher = better).
In the single-feature comparison, MI and II achieve equivalent
correlations ($\tau=0.275$ and $\tau=0.276$).
In the joint features comparison, II joint achieves $\tau=0.396$
and DII with L1 regularisation ($\lambda=0.1$) achieves
Top-10 precision~$=1.00$, while JMI greedy fails ($\tau \approx 0$)
as its sequential structure cannot identify features with zero
marginal conditional MI.
A synergistic XOR feature pair, individually uninformative but jointly
predictive, is ranked 26th and 21st by MI, 24th and 23rd by II, and
18th--19th by JMI, whereas DII and II joint rank it 1st and 2nd,
demonstrating that simultaneous optimisation is necessary to detect
pure synergistic interactions.
L1 regularisation further suppresses near-duplicate features.
Applied to real stock return data (440,402 observations, 27 indicators),
DII identifies volume and momentum features as most predictive while
suppressing 13 redundant indicators, a result undetected by per-feature
methods, confirming joint optimisation as essential when interactions
and redundancy are present.

\end{document}
